\section{Improper Integrals}

The definition of the Riemann integral was given only for bounded functions
defined on a bounded interval. We now wish to extend the definition of the
integral to unbounded functions and unbounded intervals. We will do this by
reducing, through a limit, to the case already studied.

\begin{definition}[Local Integrability]
\label{def:localmente_riemann}
Let $f\colon A \to \RR$ be a function defined on a set $A\subset \RR$. We say
that $f$ is \emph{locally Riemann-integrable} if for every closed and bounded
interval $[\alpha,\beta]\subset A$, it follows that $f$ is bounded and
Riemann-integrable on $[\alpha,\beta]$.
\end{definition}

\begin{remark}
If $f\colon A \to \RR$ is a continuous function, then on every interval
$[\alpha,\beta]\subset A$, the function $f$ is continuous and hence bounded and
Riemann-integrable by theorem~\ref{th:integrabilita_continue}. The same holds
for monotone functions: thanks to theorem~\ref{th:integrabilita_monotone}, we
know that a monotone function is bounded and Riemann-integrable on every
interval $[\alpha,\beta]$, thus it is also locally Riemann-integrable. It will
be very rare to encounter locally Riemann-integrable functions that do not fall
into these two cases.
\end{remark}

\begin{definition}[Improper Integral]
\label{def:integrale_improprio}
\index{integral!improper}
If $f$ is a locally Riemann-integrable function on the interval $[a,b)$ with
$a\in \RR$, $b\in (a,+\infty]$, we define the improper integral of $f$ on
$[a,b)$ as:
\[
  \int_a^b f(x)\, dx = \lim_{\beta \to b^-} \int_a^\beta f(x)\, dx
\]
if the limit on the right-hand side exists (finite or infinite).

If $f$ is a locally Riemann-integrable function on the interval $(a,b]$ with
$b\in \RR$, $a\in [-\infty,b)$, we define the improper integral of $f$ on
$(a,b]$ as:
\[
  \int_a^b f(x)\, dx = \lim_{\alpha \to a^+} \int_\alpha^b f(x)\, dx
\]
if the limit exists (finite or infinite).

If $f$ is a locally Riemann-integrable function on the interval $(a,b)$ with
$a\in [-\infty,+\infty)$, $b\in(a,+\infty]$, given any point $c\in (a,b)$, we
define the improper integral of $f$ on $(a,b)$ as:
\begin{align*}
  \int_a^b f(x)\, dx &=
  \int_a^c f(x)\, dx + \int_c^b f(x)\, dx \\
  &= \lim_{\alpha \to a^+}\int_\alpha^c f(x)\,dx
    + \lim_{\beta\to b^-}\int_c^\beta f(x)\, dx
\end{align*}
provided that both limits exist (finite or infinite) and are not infinite with
opposite signs (so that their sum is well-defined). According to
observation~\ref{rem:4821341}, the integral does not depend on the chosen point
$c$.

If $I$ is an interval with endpoints $a,b\in [-\infty, +\infty]$, $a<b$, and
there exist a finite number of points $x_0,\dots,x_n \in I$ such that $a = x_0 <
x_1 < \dots < x_n = b$ and if $f$ is locally Riemann-integrable on the set
$A=I\setminus\ENCLOSE{x_0, \dots, x_n}$, then we define the improper integral of
$f$ on $A$ as:
\[
  \int_a^b f(x)\, dx = \sum_{k=1}^n \int_{x_{k-1}}^{x_k} f(x)\, dx
\]
if each improper integral $\int_{x_{k-1}}^{x_k} f(x)\, dx$ exists (finite or
infinite) and if the sum is well-defined since all infinite integrals have the
same sign.

If the improper integral $\int_a^b f(x)\, dx$ exists and is finite (according to
one of the previous definitions), we say that the function $f$ is
\mymargin{integrable} \emph{integrable in the improper sense} (or
\emph{integrable in the generalized sense}) and we say that the integral
\emph{converges}. \mymargin{convergent}%
\index{convergence!integral}%
\index{integral!convergent}%
\index{integral!improper!convergent}%
If instead the integral exists but is not finite, we say that the integral
\emph{diverges}. \mymargin{divergent}%
\index{divergence!integral}%
\index{integral!divergent}%
\index{integral!improper!divergence}%
In other cases, we say that the integral is \emph{indeterminate}.
\mymargin{indeterminate}%
\index{integral!improper!indeterminate}%
\index{integral!indeterminate}%

Determining the \emph{nature} \mymargin{nature}%
\index{nature!of the integral}%
\index{integral!nature}%
\index{integral!improper!nature}%
of the integral $\int_a^b f(x)\, dx$ means stating whether such an integral is
convergent, divergent, or indeterminate.

If the limits of integration are swapped, $a>b$, it is useful to use the
convention already introduced for Riemann integrals even for improper integrals:
\[
  \int_a^b f(x) \, dx = - \int_b^a f(x)\, dx.
\]
\end{definition}

\begin{remark}\label{rem:4821341}
In the definition of the improper integral on the open interval $(a,b)$, the
choice of the point $c\in(a,b)$ does not affect the definition. Indeed, if we
choose two different points $c,c'\in(a,b)$ for every $\alpha,\beta \in (a,b)$,
we have, thanks to the additivity of the Riemann integral:
\[
  \int_\alpha^c f(x)\,dx + \int_c^\beta f(x)\,dx =
  \int_\alpha^{c'} f(x)\, dx + \int_{c'}^\beta f(x)\, dx.
\]
And thus the limits, if they exist, are equal.
\end{remark}

\begin{example}
The function $f(x)= \ln x$ can be integrated in the improper sense on the
interval $[1,+\infty)$ and we have:
\begin{align*}
 \int_1^{+\infty} f(x)\,dx
 &= \lim_{\beta\to +\infty} \int_1^\beta \ln(x)\, dx
 = \lim_{\beta\to +\infty}\Enclose{x\ln x - x}_1^\beta \\
 &= \lim_{\beta\to +\infty} \enclose{\beta \ln \beta - \beta + 1}
 = +\infty.
\end{align*}
Also on the interval $(0,1]$, the function $\ln x$ has an improper integral
\[
  \int_0^1 \ln x\, dx = \lim_{\alpha \to 0^+}\Enclose{x \ln x -x}_\alpha^1
  = \lim_{\alpha\to 0^+} (-1-\alpha \ln \alpha + \alpha)  = -1.
\]
Thus on the interval $(0,+\infty)$, the function $\ln x$ has an improper
integral
\[
  \int_0^{+\infty}\ln x\, dx = \int_0^1 \ln x\, dx + \int_1^{+\infty}\ln x\, dx
   = -1 + (+\infty) = +\infty.
\]
We will therefore say that the function $f(x)=\ln x$ is not integrable in the
improper sense on the interval $(0,+\infty)$, but the integral exists and is
divergent.
\end{example}

To abbreviate the notations, we will mean:
\[
  \Enclose {F(x)}_a^b
  = \lim_{x\to b^-}F(x) - \lim_{x\to a^+}F(x)
\]
noting that if $F$ is defined and continuous at the endpoints $a$ and $b$, this
notation coincides with the usual
\[
\Enclose{F(x)}_a^b = F(b) - F(a).
\]

We can then write:
\begin{align*}
  \int_0^1 \ln x\, dx
  &= \Enclose{x\ln x- x}_0^1
  = 1 \cdot\ln 1 -1 - \lim_{x\to 0^+} (x\ln x - x) \\
  &= -1  - 0 = -1.
\end{align*}

Note that if $f\colon(a,b) \to \RR$ is continuous and $F\colon(a,b)\to \RR$ is
its antiderivative, then we will have
\begin{align*}
  \int_a^b f(x)\, dx
   &= \lim_{\alpha\to a^+}\int_{\alpha}^c f(x)\, dx
   + \lim_{\beta\to b^-}\int_c^{\beta} f(x)\, dx \\
   & = \lim_{x\to a} (F(c)-F(x)) + \lim_{x\to b} (F(x)-F(c)) \\
   & = \lim_{x\to b} F(x) - \lim_{x\to a} F(x)
   = \Enclose{F(x)}_a^b.
\end{align*}
Thus, the fundamental theorem of calculus remains formally identical for
improper integrals.

\begin{example}
To calculate the integral
\[
  \int_{-\infty}^{+\infty} \frac{1}{1+x^2}\, dx
\]
we can choose any point $c\in \RR$ and calculate
\begin{align*}
\int_{-\infty}^{+\infty}\frac{1}{1+x^2}\, dx
&=
\lim_{\alpha\to-\infty}\int_\alpha^c \frac{1}{1+x^2}\, dx
  + \lim_{\beta\to+\infty}\int_c^\beta \frac{1}{1+x^2}\, dx\\
  &= [\arctg x]_{-\infty}^c
  + [\arctg x]_c^{+\infty}  \\
  &= \arctg c - (-\frac\pi 2) + \frac \pi 2 - \arctg c
  = \pi.
\end{align*}

But more simply, we can write:
\[
 \int_{-\infty}^{+\infty}\frac{1}{1+x^2}\, dx
 = \Enclose{\arctg x}_{-\infty}^{+\infty}
 = \frac \pi 2 - \enclose{-\frac \pi 2} = \pi.
 \]

\end{example}

\begin{remark}
If $f\colon[a,b]\to\RR$ is bounded and Riemann-integrable on $[a,b]$, then,
according to theorem~\ref{th:integrale_continuo}, we have:
\[
 \int_a^b f(x)\, dx
  = \lim_{\alpha \to a^+} \int_\alpha^b f(x)\, dx
  = \lim_{\beta\to b^-}\int_a^\beta f(x)\,dx.
\]
Thus, in this case, the improper integrals on $[a,b)$ and on $(a,b]$ coincide
with the usual Riemann integral (and are therefore convergent). Similarly, if
$f$ is locally integrable on $[a,b)$, the improper integral on $[a,b)$ and the
improper integral on $(a,b)$ coincide. The same happens if the function is
locally integrable on $(a,b)$ and we consider the set
$A=(a,b)\setminus\ENCLOSE{x_0,\dots,x_n}$. The integral on $(a,b)$ coincides
with the integral on $A$ thanks to the additivity of the integral with respect
to the domain.

This justifies using the same notation $\int_a^b$ for both the Riemann integral
on $[a,b]$ and the various improper integrals on $[a,b)$, $(a,b]$, $(a,b)$, or
$(a,b)\setminus\ENCLOSE{x_0, \dots, x_n}$.
\end{remark}

\begin{example}
The function $\sin x$, although integrable on every closed and bounded interval
(as it is a continuous function), does not admit an improper integral on the
interval $[0,+\infty)$ because the integral:
\[
  \int_0^\beta \sin(x)\, dx
  = \Enclose{-\cos(x)}_0^\beta
  = -\cos(\beta)+\cos(0) = 1-\cos(\beta)
\]
does not admit a limit as $\beta\to +\infty$.
\end{example}

\begin{example}
The function $2x/(x^2+1)$ is not integrable in the improper sense on $\RR$
because we have
\begin{align*}
\int_{0}^{+\infty} \frac {2x} {1+x^2}\, dx
 &= \Enclose{\ln (1+x^2)}_0^{+\infty} = +\infty \\
\int_{-\infty}^0 \frac{2x}{1+x^2}\, dx
 &= \Enclose{\ln(1+x^2)}_{-\infty}^0 = -\infty
\end{align*}
and $(+\infty)+ (-\infty)$ is undefined.
\end{example}

\begin{example}
We have
\begin{align*}
  \int_{-1}^2 \frac{1}{\sqrt[3]{x}}\, dx
  &= \int_{-1}^0 \frac{1}{\sqrt[3] {x}}\, dx
   + \int_{0}^2 \frac{1}{\sqrt[3] x}\, dx \\
    &= \Enclose{\frac 3 2 \sqrt[3] {x^2}}_{-1}^0  + \Enclose{\frac 3 2 \sqrt[3]
  {x^2}}_0^2
  = 0 - \frac 3 2 + \frac 3 2 \sqrt[3]{4} - 0 \\
  &= \frac 3 2 (\sqrt[3] 4-1).
\end{align*}
\end{example}

\begin{example}
We have
\[
  \int_0^{+\infty} \frac{1}{x}\, dx
  = \Enclose{\ln x}_0^{+\infty}
  = +\infty - (-\infty) = +\infty
\]
and
\[
  \int_{-\infty}^{0} \frac{1}{x}\, dx
  = \Enclose{\ln(-x)}_{-\infty}^{0}
  = -\infty - (+\infty) = -\infty.
\]
Therefore, the integral
\[
  \int_{-\infty}^{+\infty} \frac{1}{x}\, dx
\]
is not defined as it is the sum of infinities with opposite signs.
\end{example}

\begin{remark}
Be careful not to forget the "bad" points within the integration interval. If we
want to evaluate:
\[
  \int_{-1}^1 \frac{1}{x}\, dx
\]
We might be led to think that this integral is equal to:
\[
 \Enclose{\ln \abs{x}}_{-1}^1 = \ln \abs{1} - \ln \abs{-1} = 0.
\]
However, this integral is not defined because the integrand function is not
defined and is not bounded in a neighborhood of $x=0$, and therefore the
interval $[-1,1]$ must be split into the union of the two intervals $[-1,0)$ and
$(0,1]$, from which:
\[
  \int_{-1}^1 \frac{1}{x}\, dx = \Enclose{\ln \abs{x}}_{-1}^0 + \Enclose{\ln\abs{x}}_0^{1}
   = -\infty + (+\infty)
\]
and since there is a sum of infinities with opposite signs, the sum is
meaningless, and the improper integral is not defined.
\end{remark}

\begin{theorem}[Properties of the Improper Integral]
\label{th:proprieta_integrale_improprio}
If $f\le g$ and if both integrals are defined, it follows:
\index{integral!improper!monotonicity}%
\mymargin{monotonicity}%
\index{monotonicity}%
\[
  \int_a^b f(x)\, dx \le \int_a^b g(x)\, dx.
\]

Let $a,b,c \in \bar \RR$ with $a\le c \le b$. Then in the following equality
\index{integral!improper!additivity}%
\mymargin{additivity}%
\index{additivity}%
\[
  \int_a^b f(x)\, dx = \int_a^c f(x)\, dx + \int_c^b f(x)\, dx
\]
if at least one of the two sides is defined, then the other side is also
defined, and the equality holds.

Let $\lambda, \mu \in \RR$.
\index{integral!improper!linearity}%
\mymargin{linearity}%
\index{linearity}%
\[
  \int_a^b \enclose{\lambda f(x) + \mu g(x)}\, dx
  = \lambda \int_a^b f(x) + \mu \int_a^b g(x)
\]
if the right side is well-defined (the improper integrals of $f$ and $g$ exist,
and the operations make sense).

Let $g\colon (a,b)\to (c,d)$ be a $C^1$ function and $f\colon (c,d) \to \RR$ a
continuous function integrable in the improper sense on $(c,d)$. Let $-\infty
\le a < b \le +\infty$ and $-\infty \le c < d \le +\infty$ such that
\[
c = \lim_{x\to a^+} g(x),
\qquad
d = \lim_{x\to b^-} g(x).
\]
If $f$ is integrable in the improper sense on $(c,d)$, then $f \circ g$ is
integrable on $(a,b)$ and we have
\index{integral!improper!change of variable}%
\mymargin{change of variable}%
\index{change of variable}%
\[
\int_a^b f(g(x)) g'(x)\, dx = \int_c^d f(y)\, dy.
\]

If instead the limits are swapped:
\[
d = \lim_{x\to a^+} g(x),
\qquad
c = \lim_{x\to b^-} g(x)
\]
we will have, under the same hypotheses:
\[
  \int_a^b f(g(x)) g'(x) \, dx = \int_d^c f(y)\, dy.
\]

If $F$ is differentiable on $(a,b)$, $g$ is continuous, and $G$ is an
antiderivative of $g$, then the formula holds:
\index{integral!improper!by parts}%
\mymargin{integration by parts}%
\index{integration by parts}%
\[
 \int_a^b F(x)g(x)\, dx = \Enclose {F(x) G(x)}_a^b - \int_a^b F'(x)G(x)\, dx  
\]
if the quantities on the right side exist (finite or infinite) and their
difference is not an indeterminate form.
\end{theorem}
%
\begin{proof}
All these properties have already been proven for the integral of bounded
functions on bounded intervals. They can be easily extended to the lateral
improper integral by observing that the limit maintains the required properties.
Consequently, the properties hold for bilateral improper integrals by
additivity, taking care not to produce an indeterminate sum $+\infty +
(-\infty)$. Finally, by additivity, the formulas are valid for multilateral
improper integrals.
\end{proof}

\begin{exercise}\label{ex:821685}
  Calculate the following double integral:
  \[
    \int_0^{+\infty} \Enclose{\int_0^{+\infty} \sin t \cdot e^{-tx}\, dt}\, dx.
  \]
\end{exercise}
%
\begin{proof}
The innermost integral can be calculated by parts:
\begin{align*}
  I(x) &= \int_0^{+\infty} \sin t\cdot e^{-tx}\, dt \\
  &=  \Enclose{-\cos t \cdot e^{-tx}}_0^{+\infty}
    - x \int_0^{+\infty} \cos t \cdot e^{-tx}\, dt \\
  &= 0 + 1 - x\enclose{\Enclose{\sin t \cdot e^{-tx}}_0^{+\infty} 
  +x \int_0^{+\infty}\sin t \cdot e^{-tx}\, dt}
\end{align*}
from which 
\[
  I(x) = 1 - x^2 I(x), \qquad I(x) = \frac{1}{1+x^2}   
\]
and 
\[
 \int_0^{+\infty} \frac{1}{1+x^2}\, dx = \Enclose{\arctg x}_0^{+\infty} = \frac \pi 2.
\]
\end{proof}

\begin{theorem}[Integrability of Positive Functions]
\label{th:integrabilita_positive}
\mymark{**}
Let $f\colon (a,b)\to \RR$ be a locally Riemann-integrable non-negative
function: $f\ge 0$. Then the integral of $f$ exists (finite or infinite) and is
non-negative:
\[
  \int_a^b f(x)\, dx \ge 0.
\]
\end{theorem}
%
\begin{proof}
\mymark{**}
Let $c\in (a,b)$ be a fixed point and let
\[
  F(x) = \int_c^x f(t)\, dt.
\]
Since $f$ is locally Riemann-integrable, the function $F$ is well-defined, and
since $f\ge 0$, it follows that $F$ is increasing because if $x_2>x_1$, since
$f\ge 0$, we have
\[
  F(x_2)-F(x_1) = \int_{x_1}^{x_2} f(t)\, dt \ge \int_{x_1}^{x_2} 0\, dt = 0.
\]
Thus, the limits exist:
\begin{align*}
  \int_c^b f(x)\, dx  &= \lim_{\beta\to b^-} F(\beta)\\
  \int_a^c f(x)\, dx  &= \lim_{\alpha\to a^+} -F(\alpha)
\end{align*}
and since $F$ is increasing, both limits are non-negative, and therefore the sum
of the two limits is well-defined.
\end{proof}

Even if we do not know how to explicitly calculate an integral, it is often
possible to determine its convergence by comparing it, through the following
theorem, with a known integral.

\begin{theorem}[Comparison Criterion and Asymptotic Comparison]
\mymark{**}
Let $f, g \colon$ $ [a,b)\to \RR$ (with $a \le b \le +\infty$) be locally
Riemann-integrable functions. Suppose also that for every $x\in [a,b)$, we have
$f(x)\ge 0$ and $g(x)\ge 0$. Under these assumptions, we know that the two
integrals:
\begin{equation}\label{eq:921754}
  \int_a^b f(x)\, dx, \qquad \int_a^b g(x)\, dx
\end{equation}
both exist (finite or infinite) and are non-negative. Moreover, we have the
following comparison criteria.

\begin{enumerate}
\item \emph{Pointwise Comparison ``$\le$''.}
Suppose that for every $x\in [a,b)$, we have $f(x) \le g(x)$. Then we have:
\begin{equation}\label{eq:467143}
  \int_a^b f(x)\, dx \le \int_a^b g(x)\, dx.
\end{equation}
In particular, if the integral of $g$ is convergent, then the integral of $f$ is
convergent, while if the integral of $f$ is divergent, then the integral of $g$
is divergent.

\item \emph{Asymptotic Comparison ``$\ll$''.}
Suppose that for $x\to b^-$, we have $f \ll g$ (i.e., $f/g\to 0$). Then it
follows that if the integral of $g$ is convergent, then the integral of $f$ is
convergent, while if the integral of $f$ is divergent, then the integral of $g$
is divergent.

\item \emph{Asymptotic Comparison ``$\sim$''.}
Suppose that for $x\to b^-$, we have $f\sim g$ (i.e., $f/g\to 1$). Then the two
integrals~\eqref{eq:921754} have the same nature (both convergent or both
divergent).
\end{enumerate}

Analogous results hold for functions defined on an interval open on the left:
$(a, b]$ with $-\infty \le a \le b$, in which case, in the asymptotic criteria,
the limits are taken as $x\to a^+$.
\end{theorem}
%
\begin{proof}
\mymark{**}
The integrals of $f$ and $g$ exist thanks to
theorem~\ref{th:integrabilita_positive}.

If $f\le g$, the inequality~\eqref{eq:467143} is guaranteed by the monotonicity
property (theorem~\ref{th:proprieta_integrale_improprio}).

If $f(x)/g(x) \to 0$ as $x\to b^-$, it means that there exists a point
$c\in[a,b)$ such that for every $x\in [c,b)$, we have $f(x)/g(x)\le 1$, i.e.,
$f(x)\le g(x)$. Thus, by the additivity and monotonicity of the integral, we can
assert that
\[
  \int_a^b f
  = \int_a^c f + \int_c^b f
  \le \int_a^c f + \int_c^b g
  = \int_a^c f + \int_a^b g - \int_a^c g.
\]
We know that the integrals of $f$ and $g$ on the interval $[a,c]$ are convergent
since $f$ and $g$ are bounded and Riemann-integrable on $[a,c]$. Thus, if the
integral $\int_a^b g$ is convergent, then the integral $\int_a^b f$ is
convergent. Conversely, if $\int_a^b f$ is divergent, then $\int_a^b g$ is
divergent.

In the case where $f(x)/g(x)\to 1$ as $x\to b^-$, we can find a point
$c\in[a,b)$ such that $\frac 1 2 \le f(x)/g(x) \le 2$ for every $x\in [c,b)$.
Then, for every $x\in [c,b)$, we have
\[
f(x) \le 2g(x), \qquad g(x) \le 2 f(x)
\]
and we can proceed as in the previous case.
\end{proof}

In the most frequent cases, the previous theorem is applied by comparing the
function with a power:
\[
  (x-x_0)^\alpha, \qquad \alpha \in \RR.
\]

It will therefore be useful to know, as a benchmark, for which $\alpha$ these
functions have a convergent integral, as in the following example.

\begin{example}
\label{ex:416145}%
\mymark{***}%
Let $a>0$ and $p\in \RR$. The integrals
\[
  \int_{a}^{+\infty}\frac{1}{x^p}\, dx,
  \qquad
  \int_{-\infty}^{-a}\frac{1}{(-x)^p}\, dx
\]
are convergent if and only if $p>1$.

Let $x_0\in [a,b]$ and $p\in\RR$. The integrals
\[
  \int_a^{x_0} \frac{1}{(x-x_0)^p}\, dx,
  \qquad
  \int_{x_0}^b \frac{1}{(x-x_0)^p}\, dx
\]
are convergent if and only if $p<1$.

The verification is easily done by evaluating the limit of the antiderivative
$F(x) = x^{1-p}/(1-p)$ at the endpoints of the interval.
\end{example}

\begin{example}
\label{ex:integrale_gaussiana_finito}%
\index{integral!Gaussian}%
\index{function!Gaussian}%
\index{Gauss!function of}%
The integral
\[
 \int_{-\infty}^{+\infty}e^{-x^2}\, dx
\]
is convergent because for $x\to +\infty$, it holds that $e^{-x^2} \ll 1/x^2$,
and the integral of $1/x^2$ is convergent (example~\ref{ex:416145}) in a
neighborhood of $+\infty$ and $-\infty$. In
exercise~\ref{ex:integrale_gaussiana}, we will demonstrate that this integral
equals $\sqrt \pi$.
\end{example}

\begin{example}[Euler's Gamma Function]%
  \label{ex:gamma_eulero}%
  \index{$\Gamma$ Euler's function}%
  \index{function!$\Gamma$ of Euler}%
  We define $\Gamma\colon (0,+\infty) \to \RR$ as
  \[
    \Gamma(x) = \int_0^{+\infty} e^{-t} t^{x-1}\, dt.
  \]
    The integral converges for every $x>0$ because, setting $f(t) = e^{-t}
  t^{x-1}$, for $t\to 0^+$, we have $f(t) \sim t^{x-1}$, which has a convergent
  integral in a neighborhood of $0$, while for $t\to +\infty$, we have $f(t) \ll
  e^{-t/2}$, which has a convergent integral in a neighborhood of $+\infty$.
  
  Integrating by parts, we obtain an interesting property of the Gamma function:
  \begin{align*}
  \int_0^{+\infty} e^{-t}t^{x}\, dt
  &= \Enclose{-e^{-t}t^{x}}_0^{+\infty}
  + \int_0^{+\infty} e^{-t}xt^{x-1}\, dt\\
  &= x\int_0^{+\infty} e^{-t}t^{x-1}\, dt
  \end{align*}
  that is, $\Gamma(x+1) = x\Gamma(x)$.
    Observing that $\Gamma(1)=1=0!$, we can thus prove, by induction, that
  $\Gamma(n+1) = n!$ for every $n\in \NN$. We have therefore found a function
  that extends the factorial from $\NN$ to the entire interval of real numbers
  $(-1,+\infty)$.
  
    In chapter~\ref{sec:scambio_integrale_limite}, we will demonstrate that the
  Gamma function is differentiable.
  
    We are not currently able to calculate the value of $\Gamma(x)$ if $x$ is
  not an integer, but we can relate the value of $\Gamma(1/2)$ to the integral
  of the Gaussian function by making the change of variable $t=s^2$, $dt=2s ds$:
  \begin{align*}
  \Gamma(1/2) 
  &= \int_0^{+\infty} \frac{e^{-t}}{\sqrt t}\, dt  
  = \int_0^{+\infty} \frac{e^{-s^2}}{s} \cdot 2s\, ds
  = \int_{-\infty}^{+\infty} e^{-s^2}\, ds.
  \end{align*}
    In exercise~\ref{ex:integrale_gaussiana}, we will find the value of the
  integral of the Gaussian function and will be able to conclude that
  $\Gamma(1/2) = \sqrt \pi$. This allows us to calculate the value of $\Gamma$
  on all half-integers:
  \[
  \Gamma(3/2) = \frac 1 2 \Gamma(1/2) = \frac{\sqrt \pi}{2}, \qquad
  \Gamma(5/2) = \frac 3 2 \Gamma(3/2) = \frac{3\sqrt \pi}{4}, \dots 
  \]
\end{example}
    
\begin{example}
The integral
\[
  \int_{-\infty}^{+\infty}\frac{1}{1+x^2}\, dx
\]
is convergent because for $x\to +\infty$, it holds that $1/(1+x^2) \sim 1/x^2$.

The integral
\[
  \int_1^{+\infty} \frac{1}{\ln x}\, dx
\]
is divergent because for $x\to +\infty$, we have $1/\ln x \gg 1/x$, and for
$x\to 1^+$, we have $1/\ln x \sim 1/(x-1)$. For $p=1$, the integrals of $1/x$ at
$+\infty$ and of $1/(x-1)$ at $1$ are both divergent.

If $f\colon [a,+\infty)$ is a locally Riemann-integrable function and if
$f(x)\to \ell \neq 0$ as $x\to +\infty$, then the integral of $f$ is divergent.
If $\ell>0$, there will exist $c>a$ such that for $x>c$, we have $f(x)\ge 0$
(persistence of sign). Since $f(x)\sim \ell$ as $x\to +\infty$, we can conclude
that $\int_c^{+\infty} f(x) = +\infty$ and therefore also $\int_a^{+\infty} f(x)
= +\infty$. If $\ell<0$, we can change the sign of $f$ and repeat the previous
argument, thus discovering that in this case, the integral of $f$ is $-\infty$.
\end{example}

\begin{exercise}
Show with an example that if a function $f$ has a convergent integral on the
interval $[0,+\infty)$, it is not necessarily the case that $f(x)\to 0$ as $x\to
+\infty$.
\end{exercise}

% \begin{exercise}[difficult]
% Show that the integral
% \[
%   \int_0^{+\infty} \exp\enclose{x^2 \ln \frac{2+\cos(2\pi x)}{3}} \, dx
% \]
% is finite. But the integrand function does not tend to zero as $x\to +\infty$.
% \end{exercise}

\begin{exercise}
Determine whether the following integrals are convergent:
\[
  \int_0^{+\infty} \frac{1}{(x^2-1)\cdot \ln^2 x}\, dx, \qquad
  \int_0^{+\infty} \frac{\ln^2 x}{x^2-1}\, dx.  
\]
\end{exercise}

\begin{exercise}
Determine the values of $p\in \RR$ for which the following integral is
convergent:
  \[
  \int_0^{+\infty} \frac{\sqrt x-\ln x}{(x-\sin x)^p}\, dx.  
  \]
\end{exercise}


\begin{theorem}[Criterion of Absolute Convergence]
\label{th:convergenza_assoluta_integrale}
\mymark{**}%
Let $f\colon (a,b)\to \RR$ be a locally Riemann-integrable function. Then
$\abs{f}$ is also locally Riemann-integrable, and if
\[
\int_a^b \abs{f(x)}\, dx < +\infty
\]
(i.e., $\abs{f}$ is integrable in the improper sense), then $f$ is also
integrable in the improper sense, and
\[
  \abs{\int_a^b f(x)\,dx}
  \le \int_a^b \abs{f(x)}\, dx < +\infty.
\]
\end{theorem}
%
%
\begin{proof}
Let $f = f^+ - f^-$ with $f^+\ge 0$ and $f^-\ge 0$. If $f$ is locally
Riemann-integrable, then $f^+$ and $f^-$ are locally Riemann-integrable (thanks
to theorem~\ref{th:reticolo}). Note that we have $f = f^+ - f^-$ and $\abs{f} =
f^+ + f^-$. Since $0\le f^+ \le \abs{f}$, we can apply the comparison criteria
and assert that $f^+$ has a convergent integral. The same holds for $f^-$ and
thus for $f=f^+ - f^-$. Moreover,
\[
  \abs{\int_a^b f }
  = \abs{\int_a^b f^+ - \int_a^b f^-}
  \le \int_a^b f^+ + \int_a^b f^-
  = \int_a^b \abs{f}.
\]
\end{proof}

\begin{definition}[Absolute Convergence]
\index{integral!improper!absolute convergence}%
\index{integral!absolute convergence}%
\index{convergence!absolute!integral}%
\index{integrable!absolutely}%
\index{absolutely!integrable}%
When
\[
  \int_a^b \abs{f(x)}\, dx <+\infty
\]
we say that the integral of $f$ is \emph{absolutely convergent} or that $f$ is
\emph{absolutely integrable} (in the improper sense). The previous theorem
guarantees that an absolutely integrable function (if locally
Riemann-integrable) is integrable.
\end{definition}

\begin{example}
The integral
\[
  \int_0^{+\infty} \sin(x^2)\cdot e^{-x}\, dx
\]
is absolutely convergent (and therefore convergent) because the integrand
function is continuous, and we have
\[
 \abs{\sin(x^2)\cdot e^{-x}} \le e^{-x}
\]
from which, by comparison,
\[
 \int \abs{\sin(x^2)\cdot e^{-x}}\, dx  \le
  \int_0^{+\infty} e^{-x}\, dx < +\infty.
\]
\end{example}

\begin{example}[Function Integrable but Not Absolutely]
\label{ex:48864}%
\mymark{*}%
We want to show that the integral
\[
  \int_\pi^{+\infty} \frac{\sin x}{x}\, dx  
\]
is convergent but not absolutely convergent.

One way to do this is to relate it to series. Let 
\[
  a_k = \int_{k\pi}^{(k+1)\pi} \frac{\sin x}{x}\, dx
\]
we have
\[
  \abs{a_k} = \int_{k\pi}^{(k+1)\pi} \frac{\abs{\sin x}}{x}\, dx
\]
and thus 
\[
\int_\pi^{+\infty} \frac{\sin x}{x}\, dx
= \sum_{k=1}^{+\infty} a_k, \qquad
\int_\pi^{n\pi} \frac{\abs{\sin x}}{x}\, dx
= \sum_{k=1}^{n-1} \abs{a_k}.
\]



Note that 
\[
   \int_\pi^{n\pi} \frac{\sin x}{x}\, dx
   = \sum_{k=1}^{n-1} \int_{k\pi}^{(k+1)\pi} \frac{\sin x}{x}\, dx
\]
and that 



Integrating by parts,
\begin{align*}
  \int_1^{+\infty} \frac{\sin x}{x}\, dx
  &= \Enclose{\frac{-\cos x}{x}}_1^{+\infty} -
  \int_1^{+\infty} \frac{\cos x}{x^2} \, dx \\
  &= \cos 1 - \int_1^{+\infty} \frac{\cos x}{x^2}\, dx.
\end{align*}
Now observe that the function $\frac{\cos x}{x^2}$ is integrable by the
criterion of absolute convergence, since:
\[
  \abs{\frac{\cos x}{x^2}} \le \frac{1}{x^2}
\]
which is integrable on $[1,+\infty)$. Thus, our function has a convergent
integral also on $[1,+\infty)$ and therefore has a convergent integral on the
entire $(0,+\infty)$. In exercise~\ref{ex:integrale_sin_integrale}, we will see
that $\int_0^{+\infty} \frac{\sin x}{x} = \frac \pi 2$.

However, this function is not absolutely integrable. Indeed, for every $k\in
\NN$, we have
\begin{align*}
  \int_{k\pi}^{(k+1)\pi} \abs{\frac{\sin x} x}\, dx
  &\ge \int_{k\pi}^{(k+1)\pi} \frac{\abs{\sin x}}{(k+1)\pi}\, dx\\
  &= \frac{\int_0^\pi \sin x \, dx}{(k+1)\pi}
  = \frac{2}{(k+1)\pi}
\end{align*}
from which
\begin{align*}
  \int_0^{+\infty} \abs{\frac{\sin x}{x}}\, dx
  & = \lim_{b\to +\infty} \int_0^b \abs{\frac{\sin x}{x}}\, dx 
  \ge \lim_{b\to+\infty} \sum_{k=0}^{\lfloor \frac{b}{\pi}\rfloor-1} \int_{k\pi}^{(k+1)\pi}\abs{\frac{\sin x} x}\, dx\\
  &= \sum_{k=0}^{+\infty} \frac{2}{(k+1)\pi} = +\infty.
\end{align*}
\end{example}

\begin{exercise}
    Try to apply the method used in example \ref{ex:48864} to a generic integral
  of the form $\int_1^{+\infty} f\cdot g$, identifying the hypotheses on $f$ and
  $g$ that guarantee the convergence of the integral. You should obtain a
  criterion analogous to Dirichlet's theorem~\ref{th:dirichlet} for the
  convergence of series.
\end{exercise}


The criteria we have seen so far highlight a remarkable analogy between series
and integrals. Determining the convergence of an integral is, in general,
simpler than determining the convergence of the corresponding series. In the
following exercise, for example, if $\alpha=1$, it is sufficient to find an
antiderivative of the integrand functions to determine the nature of the
integral.

\begin{exercise}
Determine for which $\alpha,\beta\in \RR$ the integrals are convergent:
\[
  \int_0^{\frac 1 2} \frac{1}{x^\alpha \abs{\ln x}^\beta}\, dx,
  \qquad
  \int_2^{+\infty} \frac{1}{x^\alpha \ln^\beta x} \, dx.
\]
\end{exercise}

\begin{theorem}[Connection Between Series and Improper Integrals]
\mymark{**}
Let $f\colon$ $[1,+\infty)\to \RR$ be a non-negative decreasing function, and
let $a_k=f(k)$. Then the series
\[
   \sum_{k=1}^{+\infty} a_k
\]
is convergent if and only if the integral
\[
  \int_1^{+\infty} f(x)\, dx.
\]
is convergent.

And, more precisely, for every $n\in \NN\cup\ENCLOSE{+\infty}$, we have
\begin{equation}\label{eq:39185}
  0
  \le \sum_{k=1}^{n} a_k - \int_1^{n}f(x)\, dx
  \le a_1.
\end{equation}

\end{theorem}
%
\begin{proof}
\mymark{**}
Since the function $f$ is non-negative and monotonic, the improper integral of
$f$ on $[1,+\infty)$ exists (finite or infinite) by
theorem~\ref{th:integrabilita_positive}. The series $\sum f(k)$, being
non-negative, is also determined.

Since $f$ is decreasing, for every $x \in [k,k+1]$, we have
\[
  f(k+1) \le f(x) \le f(k)
\]
integrating over $[k,k+1]$, we obtain
\[
  f(k+1) \le \int_{k}^{k+1} f(x)\, dx \le f(k)
\]
and summing for $k=1,\dots, n-1$, we obtain:
\[
  \sum_{k=2}^{n} f(k) \le \int_{1}^{n} f(x)\, dx \le \sum_{k=1}^{n-1} f(k)
\]
which implies~\eqref{eq:39185} if $n\in \NN$. Letting
$n\to+\infty$,~\eqref{eq:39185} remains valid. The first inequality allows us to
estimate the series with the integral, and the second allows us to estimate the
integral with the series. Thus, if one of the two converges, the other also
converges, as we wanted to prove.
\end{proof}

\begin{example}[Asymptotic Estimate of $\ln n!$]
\label{ex:498124}%
\index{factorial!asymptotic estimate}%
\mymargin{asymptotic estimate $\ln n!" $}%
Note that
\[
  \ln (n!) = \ln \prod_{k=1}^n k = \sum_{k=1}^n \ln k.
\]
We apply the same procedure used in the previous theorem to the function $f(x) =
\ln x$. Since the logarithm is increasing (instead of decreasing), we will
obtain reversed but analogous estimates. Repeating the calculations carefully,
we obtain:
\[
    \int_1^n \ln(x)\, dx \le  \sum_{k=2}^n \ln k \le \int_2^{n+1} \ln(x) \, dx
\]
and recalling that $\int \ln x = x \ln x - x$, we obtain
\[
  n \ln n - n + 1 \le \ln(n!) \le (n+1) \ln (n+1) - n - 2\ln 2 + 1
\]
from which, dividing both sides by $n \ln n$ and letting $n\to +\infty$, we find
\[
 \frac{\ln n!}{n \ln n}\to 1
\]
that is,
\[
  \ln (n!) \sim n \ln n \qquad \text{as $n\to +\infty$.}
\]

Note that, in general, we cannot take the exponential of both sides of an
asymptotic equivalence and hope that the asymptotic equivalence remains valid.
In fact, finding an asymptotic estimate for $n!$ is much more complicated
(compared to the estimate we found for its logarithm). We will do this in
theorem~\ref{th:stirling} by finding Stirling's formula. See also
exercise~\ref{ex:7340098}.

\end{example}